% year: תשע"ז
% type: מבחן
% semester: א
% moed: א
% number: 1
% points: 20
% categories: taylor
% source: exams/מבחנים/תשעז/מועד א תשעז.pdf

\begin{question}
השתמשו בנוסחת טיילור כדי לחשב $\sqrt[3]{129}$ עם שגיאה לא גדולה מ-$10^{-3}$.
הסבירו כיצד הערכתם את השגיאה.
\end{question}

\begin{hint}
$129$ קרוב ל-$125=5^3$. פתחו את $f(x)=\sqrt[3]{x}$ סביב $x_0=125$.
\end{hint}

\begin{hint}
העריכו את השארית בצורת לגרנז' $R_n(x)=\frac{f^{(n+1)}(c)}{(n+1)!}(x-x_0)^{n+1}$ עם $125<c<129$, ובדקו איזה $n$ מספיק.
\end{hint}

\begin{solution}
נגדיר $f(x)=\sqrt[3]{x}=x^{1/3}$ ונפתח סביב $x_0=125$ (כי $\sqrt[3]{125}=5$ ידוע במדויק), עם $x=129$, כלומר $x-x_0=4$.

\textbf{נגזרות:}
\[
f'(x)=\frac13 x^{-2/3},\qquad f''(x)=-\frac29 x^{-5/3}.
\]
לכן $f(125)=5$, $f'(125)=\frac13\cdot\frac{1}{25}=\frac{1}{75}$.

\textbf{פולינום טיילור מסדר 1:}
\[
P_1(129)=f(125)+f'(125)\cdot 4=5+\frac{4}{75}=\frac{379}{75}\approx 5.05333.
\]

\textbf{הערכת השגיאה.} לפי משפט טיילור עם שארית לגרנז', קיימת נקודה $c$ בין $125$ ל-$129$ כך ש-
\[
R_1(129)=\frac{f''(c)}{2!}\cdot 4^2=-\frac{2}{9}c^{-5/3}\cdot\frac{16}{2}=-\frac{16}{9\,c^{5/3}}.
\]
מכיוון ש-$c>125$, מתקיים $c^{5/3}>125^{5/3}=5^5=3125$, ולכן
\[
|R_1(129)|<\frac{16}{9\cdot 3125}=\frac{16}{28125}\approx 5.7\cdot 10^{-4}<10^{-3}.
\]
לכן פולינום טיילור מסדר 1 מספיק.

\textbf{תשובה:}
\[
\sqrt[3]{129}\approx 5+\frac{4}{75}=\frac{379}{75}\approx 5.0533,
\]
עם שגיאה קטנה מ-$5.7\cdot10^{-4}<10^{-3}$. (השארית שלילית, כלומר הקירוב גדול מהערך האמיתי; ואכן $\sqrt[3]{129}\approx 5.05277$.)
\end{solution}
